\subsection{Generation of a Riesz space by its positive elements}

Let \(\mathbf{E}\) be an ordered vector space; the set \(\mathrm{P}\) of elements \(\geqslant 0\) of \(\mathbf{E}\) is a convex cone with vertex 0, that is (TVS, II, §2, No.~4), a set such that

\(P + P \subset P\) and \(\lambda P \subset P\) for all \(\lambda > 0\) . Conversely, if, in a vector space E over R, P is a convex cone with vertex 0, such that \(P \cap (-P) = \{0\}\) (in other words, a pointed and proper convex cone), one knows (loc. cit.) that the relation \(y - x \in P\) is an order relation (denoted \(x \leqslant y\) ) compatible with the vector space structure of E. For this order structure to define a Riesz space structure on E, it is necessary and sufficient that:

\(1^{\circ}\) P generates E, that is, every \(z \in \mathbf{E}\) is of the form \(y - x\) , where \(x\) and \(y\) belong to P;

\(2^{\circ}\) P satisfies one of the following two conditions:

\begin{enumerate}
\def\labelenumi{\alph{enumi})}
\item
  every pair of elements of \(\mathbf{P}\) has a supremum in \(\mathbf{P}\) ;
\item
  every pair of elements of P has an infimum in P (A, VI, §1, No.~9, Prop. 8).
\end{enumerate}

